\documentclass{beamer}

%lembre-se de configurar o pdflatex com:
%pdflatex -shell-escape -synctex=1 -interaction=nonstopmode  -output-directory=build  %.tex

\usepackage{aula}

\title{Introdução a Algoritmos}
\date{\today}

\begin{document}
    
\input{capa}



    %------------------------------------------------
    
    \begin{frame}{Regressão Linear}
        
        Modelo de regressão linear:
        
        \[
        y = w^T x + b
        \]
        
        onde
        
        \begin{itemize}
            \item $x$ vetor de entrada
            \item $w$ pesos
            \item $b$ bias
        \end{itemize}
        
        Objetivo: encontrar $w$ e $b$ que minimizam o erro de previsão.
        
    \end{frame}
    
    %------------------------------------------------
    
    \begin{frame}{Representação Gráfica}
        
        \centering
        
        \begin{tikzpicture}[>=Stealth]
            
            \node (x1) at (0,1) {$x_1$};
            \node (x2) at (0,0) {$x_2$};
            \node (x3) at (0,-1) {$x_3$};
            
            \node[circle,draw] (sum) at (3,0) {$\sum$};
            
            \node (y) at (5,0) {$y$};
            
            \draw[->] (x1) -- node[above] {$w_1$} (sum);
            \draw[->] (x2) -- node[above] {$w_2$} (sum);
            \draw[->] (x3) -- node[below] {$w_3$} (sum);
            
            \draw[->] (sum) -- (y);
            
            \node at (3,-1.5) {$+\,b$};
            
        \end{tikzpicture}
        
    \end{frame}
    
    %------------------------------------------------
    
    \begin{frame}{Neurônio Artificial}
        
        Um neurônio calcula
        
        \[
        z = w^T x + b
        \]
        
        e aplica uma função de ativação
        
        \[
        y = \sigma(z)
        \]
        
        onde $\sigma$ pode ser:
        
        \begin{itemize}
            \item ReLU
            \item sigmoid
            \item tanh
        \end{itemize}
        
    \end{frame}
    
    %------------------------------------------------
    
    \begin{frame}{Diagrama de um Neurônio}
        
        \centering
        
        \begin{tikzpicture}[>=Stealth]
            
            \node (x1) at (0,1) {$x_1$};
            \node (x2) at (0,0) {$x_2$};
            \node (x3) at (0,-1) {$x_3$};
            
            \node[circle,draw,minimum size=1.5cm] (neuron) at (3,0) {$\sigma$};
            
            \node (y) at (5,0) {$y$};
            
            \draw[->] (x1) -- node[above] {$w_1$} (neuron);
            \draw[->] (x2) -- node[above] {$w_2$} (neuron);
            \draw[->] (x3) -- node[below] {$w_3$} (neuron);
            
            \draw[->] (neuron) -- (y);
            
            \node at (3,-1.5) {$+\,b$};
            
        \end{tikzpicture}
        
    \end{frame}
    
    %------------------------------------------------
    
    \begin{frame}{Equivalência}
        
        Se removemos a função de ativação:
        
        \[
        \sigma(z) = z
        \]
        
        então
        
        \[
        y = w^T x + b
        \]
        
        que é exatamente o modelo de regressão linear.
        
    \end{frame}
    
    %------------------------------------------------
    
    \begin{frame}{Comparação Estrutural}
        
        \begin{center}
            
            \begin{tikzpicture}[>=Stealth]
                
                \node at (0,2) {\textbf{Regressão Linear}};
                
                \node (x1) at (-2,1) {$x_1$};
                \node (x2) at (-2,0) {$x_2$};
                \node (x3) at (-2,-1) {$x_3$};
                
                \node[circle,draw] (sum) at (0,0) {$\sum$};
                
                \node (y) at (2,0) {$y$};
                
                \draw[->] (x1) -- (sum);
                \draw[->] (x2) -- (sum);
                \draw[->] (x3) -- (sum);
                
                \draw[->] (sum) -- (y);
                
            \end{tikzpicture}
            
        \end{center}
        
        \vspace{0.5cm}
        
        \begin{itemize}
            \item rede neural com 1 neurônio
            \item sem ativação
        \end{itemize}
        
    \end{frame}
    
    %------------------------------------------------
    
    \begin{frame}{Rede Neural Profunda}
        
        \centering
        
        \begin{tikzpicture}[>=Stealth]
            
            \node[circle,draw] (x1) at (0,1) {};
            \node[circle,draw] (x2) at (0,0) {};
            \node[circle,draw] (x3) at (0,-1) {};
            
            \node[circle,draw] (h1) at (2,1) {};
            \node[circle,draw] (h2) at (2,0) {};
            \node[circle,draw] (h3) at (2,-1) {};
            
            \node[circle,draw] (y) at (4,0) {};
            
            \draw[->] (x1) -- (h1);
            \draw[->] (x1) -- (h2);
            \draw[->] (x2) -- (h2);
            \draw[->] (x3) -- (h3);
            
            \draw[->] (h1) -- (y);
            \draw[->] (h2) -- (y);
            \draw[->] (h3) -- (y);
            
        \end{tikzpicture}
        
    \end{frame}
    
    %------------------------------------------------
    
    \begin{frame}{Conclusão}
        
        \begin{itemize}
            
            \item Regressão linear pode ser vista como um neurônio artificial.
            
            \item Redes neurais são composições de transformações lineares e não lineares.
            
            \item Deep learning generaliza modelos lineares clássicos.
            
        \end{itemize}
        
    \end{frame}
   
\begin{frame}{Erro de Previsão}
    Para cada ponto temos um erro:
    
    \[
    e_i = y_i - \hat{y}_i
    \]
    
    onde
    
    \[
    \hat{y}_i = wx_i + b
    \]
    
    Erro total precisa agregar os erros de todos os pontos.
\end{frame}

%------------------------------------------------

\begin{frame}{Método dos Mínimos Quadrados}
    A abordagem clássica é minimizar a soma dos erros ao quadrado.
    
    Função de custo:
    
    \[
    J(w,b) =
    \sum_{i=1}^{n}
    (y_i - (wx_i + b))^2
    \]
    
    Objetivo:
    
    \[
    \min_{w,b} J(w,b)
    \]
    
    Quadrados são usados para:
    
    \begin{itemize}
        \item evitar cancelamento de erros positivos e negativos
        \item penalizar erros grandes
    \end{itemize}
\end{frame}

%------------------------------------------------

\begin{frame}{Forma Matricial}
    Podemos escrever o modelo linear como:
    
    \[
    y = X\beta
    \]
    
    onde
    
    \[
    \beta =
    \begin{bmatrix}
        b \\
        w
    \end{bmatrix}
    \]
    
    A solução de mínimos quadrados é
    
    \[
    \beta = (X^T X)^{-1} X^T y
    \]
    
    Essa é uma solução analítica.
\end{frame}

%------------------------------------------------

\begin{frame}{Otimização por Gradiente}
    Nem sempre temos solução analítica.
    
    Podemos usar descida do gradiente:
    
    \[
    \theta_{t+1}
    =
    \theta_t
    -
    \eta \nabla J(\theta)
    \]
    
    onde
    
    \begin{itemize}
        \item $\theta$ = parâmetros
        \item $\eta$ = taxa de aprendizado
        \item $\nabla J$ = gradiente da função de erro
    \end{itemize}
\end{frame}

%------------------------------------------------

\begin{frame}{Redes Neurais}
    Uma rede neural pode ser vista como uma composição de funções.
    
    Exemplo simples:
    
    \[
    \hat{y} = f(wx + b)
    \]
    
    com função de ativação $f$.
    
    O objetivo continua sendo minimizar uma função de erro:
    
    \[
    J(y,\hat{y})
    \]
\end{frame}

%------------------------------------------------

\begin{frame}{Erro Quadrático Médio}
    Uma função de perda comum é o erro quadrático médio.
    
    \[
    J = MSE = 
    \sum (y - \hat{y})^2
    \]
    
    Observe que esta é exatamente a mesma ideia
    do método de mínimos quadrados.
\end{frame}

%------------------------------------------------

\begin{frame}{Backpropagation}
    Backpropagation é um algoritmo para calcular gradientes em redes neurais.
    
    Ele usa a regra da cadeia para calcular
    
    \[
    \frac{\partial J}{\partial w}
    \]
    
    para todos os pesos da rede.
    
    Depois disso, aplicamos gradiente descendente para atualizar os pesos.
\end{frame}

%------------------------------------------------

\begin{frame}{Regra da Cadeia}
    Considere
    
    \[
    J = (y - \hat{y})^2
    \]
    
    e
    
    \[
    \hat{y} = f(wx+b)
    \]
    
    O gradiente é calculado como
    
    \[
    \frac{\partial J}{\partial w}
    =
    \frac{\partial J}{\partial \hat{y}}
    \cdot
    \frac{\partial \hat{y}}{\partial w}
    \]
    
    Backpropagation aplica essa ideia camada por camada.
\end{frame}

%------------------------------------------------

\begin{frame}{Relação entre os Conceitos}
    
    \begin{block}{Mínimos Quadrados}
        Define a função de erro a ser minimizada.
    \end{block}
    
    \begin{block}{Gradiente Descendente}
        Método para minimizar a função de erro.
    \end{block}
    
    \begin{block}{Backpropagation}
        Algoritmo eficiente para calcular gradientes em redes profundas.
    \end{block}
    
\end{frame}

%------------------------------------------------

\begin{frame}{Resumo}
    
    \begin{itemize}
        \item Regressão linear modela relações lineares entre variáveis.
        \item Mínimos quadrados define uma função de erro baseada em erro quadrático.
        \item Gradiente descendente pode ser usado para minimizar essa função.
        \item Backpropagation calcula gradientes em redes neurais usando regra da cadeia.
        \item Redes neurais frequentemente usam erro quadrático como função de perda.
    \end{itemize}
    
\end{frame}

\end{document}
